% Part: computability
% Chapter: recursive-functions
% Section: normal-form

\documentclass[../../../include/open-logic-section]{subfiles}

\begin{document}

\olfileid{cmp}{rec}{nft}
\olsection{నియత రూప సిద్ధాంతం}

\begin{thm}[క్లీనీ నియత రూప సిద్ధాంతం]
\ollabel{thm:kleene-nf}
కింది ధర్మం గల ఆదిమ పునరావృత్త సంబంధం $T(e,x,s)$, ఆదిమ పునరావృత్త
ప్రమేయం $U(s)$ ఉన్నాయి: $f$ ఏదైనా పాక్షిక పునరావృత్త ప్రమేయమైతే,
ఏదో~$e$కు ప్రతి $x$కూ
\[
f(x) \simeq U(\umin{s}{T(e, x, s)})
\]
అవుతుంది.
\end{thm}

\begin{explain}
నియత రూప సిద్ధాంతపు నిరూపణ విస్తృతమైనది, కానీ మూల భావం సరళం. ప్రతి
పాక్షిక పునరావృత్త ప్రమేయానికి ఒక \emph{సూచిక}~$e$ ఉంటుంది; అంతర్బోధతో,
అది ఆ ప్రమేయపు ప్రోగ్రాము లేదా నిర్వచనాన్ని సంకేతీకరించే సంఖ్య.
$f(x)\fdefined$ అయితే, గణనను క్రమబద్ధంగా లిఖించి ఏదో సంఖ్య~$s$తో
సంకేతీకరించవచ్చు. $s$ అనేది ఇన్‌పుట్~$x$పై $f$ గణనను
సంకేతీకరిస్తుందనే వాస్తవాన్ని $x$, నిర్వచనం~$e$లను మాత్రమే ఉపయోగించి
ఆదిమ పునరావృత్తంగా పరీక్షించవచ్చు. అందువల్ల ``సూచిక~$e$ గల ప్రమేయానికి
ఇన్‌పుట్~$x$పై ఒక గణన ఉంది; $s$ ఆ గణనకు సంకేతం'' అనే సంబంధం~$T$
ఆదిమ పునరావృత్త సంబంధం. పూర్తి గణనా లేఖనం~$s$ ఇచ్చినప్పుడు, $s$ యొక్క
ఫలితమైన $f(x)$ విలువను కూడా $s$ నుంచి ఆదిమ పునరావృత్తంగా పొందవచ్చు.

ఏ పాక్షిక పునరావృత్త ప్రమేయపు నిర్వచనానికైనా ఒకే అపరిమిత అన్వేషణ
చాలని నియత రూప సిద్ధాంతం చూపుతుంది. ప్రాథమికంగా, సూచిక~$e$ గల
ప్రమేయానికి ఇన్‌పుట్~$x$పై గణనను సంకేతీకరించే సంఖ్య దొరికేవరకు అన్ని
సంఖ్యల్లో వెతకవచ్చు. సంఖ్యలు~$e$ను పాక్షిక పునరావృత్త ప్రమేయాల
``పేర్లు''గా వాడి, సిద్ధాంతంలోని సమీకరణ నిర్వచించే ప్రమేయం~$f$కు
$\cfind{e}$ అని రాయవచ్చు. ఏ పాక్షిక పునరావృత్త ప్రమేయానికైనా ఒకటికన్నా
ఎక్కువ సూచికలు ఉండవచ్చని గమనించండి; నిజానికి ప్రతి పాక్షిక పునరావృత్త
ప్రమేయానికి అనంతంగా అనేక సూచికలు ఉంటాయి.
\end{explain}
\end{document}
